\begin{table}[H]
\centering
\caption{Relative performance comparison by MEAN_SPL using the Friedman test with post-hoc Nemenyi analysis.}
\label{tab:stats_mean_spl}
\begin{threeparttable}
\begin{tabular}{lccccc}
\toprule
Model & Mean MEAN_SPL & Std. & Avg. rank & Sig. wins & Sig. losses \\
\midrule
DeepSets\_Q & 0.2139 & 0.0611 & 3.289 & 6 & 0 \\
M1\_AggHistFutureSummary & 0.2106 & 0.0526 & 3.364 & 6 & 0 \\
M3\_FullSkuTemporalCNN\_Q & 0.2102 & 0.0495 & 3.612 & 6 & 0 \\
SetTransformer\_Q & 0.2210 & 0.0719 & 4.074 & 6 & 0 \\
M0\_AggHistOnly & 0.2199 & 0.0659 & 4.264 & 6 & 0 \\
BottomUpGlobalHistGB & 0.2618 & 0.1240 & 6.306 & 3 & 5 \\
AggregateHistGB & 0.2606 & 0.0957 & 7.174 & 2 & 5 \\
ChildSummaryHistGB & 0.2595 & 0.0918 & 7.198 & 2 & 5 \\
AggregateElasticNet & 0.3405 & 0.2365 & 8.149 & 0 & 6 \\
ChildSummaryElasticNet & 0.3649 & 0.2661 & 9.157 & 0 & 8 \\
SeasonalNaive & 0.3391 & 0.1743 & 9.413 & 0 & 8 \\
\bottomrule
\end{tabular}
\begin{tablenotes}
\footnotesize \item The Friedman test uses 121 aligned series-level blocks (task,agg\_id). The test statistic is $\chi^2=614.5049$ with $p=1.372e-125$. Average ranks are computed with lower MEAN_SPL indicating better performance. Nemenyi significant wins/losses are counted at $\alpha=0.10$ using critical difference $CD=1.2697$.
\end{tablenotes}
\end{threeparttable}
\end{table}
